\unit{Exposición 149 — Teoremas de dualidad para los haces algebraicos coherentes}
\sid{149}
\src{169}
\seminar{(Mayo de 1957)}
\paperhead{TEOREMAS DE DUALIDAD PARA LOS HACES ALGEBRAICOS COHERENTES}%
  {por Alexandre GROTHENDIECK}

Los resultados que siguen, inspirados en el \fquote{teorema de dualidad algebraica}
de Serre, fueron obtenidos en los inviernos de 1955 y 1956. Se demuestran muy
sencillamente con ayuda de resultados bastante elementales sobre la cohomología de los espacios
proyectivos [3] y mediante el uso intensivo del álgebra homológica de
Cartan--Eilenberg, en la forma de [2].

\fsection{1}{Los Ext de haces de módulos ({[2]}, caps. 3 y 4)}

Sea $X$ un espacio topológico provisto de un haz $\Osh$ de anillos con unidad
(no necesariamente conmutativos). Consideramos la categoría abeliana
$\Ccat^{\Osh}$ de los haces de $\Osh$-módulos, llamados también $\Osh$-Módulos.
Se sabe que todo objeto de esta categoría admite una resolución inyectiva, lo que
permite definir en ella los funtores $\operatorname{Ext}$ con las propiedades
formales bien conocidas. Más precisamente, para evitar confusiones,
denotamos por $\operatorname{Hom}_{\Osh}(X;A,B)$, o simplemente por
$\operatorname{Hom}(X;A,B)$, el grupo abeliano de los $\Osh$-homomorfismos de $A$
en $B$, mientras que $\Homsh_{\Osh}(A,B)$ denotará el haz de gérmenes
de homomorfismos de $A$ en $B$ ($A,B\in\Ccat^{\Osh}$). Para
$A\in\Ccat^{\Osh}$ fijo se definen funtores $h_A$ y $\hsh_A$ con valores, respectivamente,
en la categoría $\Ccat$ de los grupos abelianos y en la categoría
$\Ccat^{\Zsh}$ de los haces abelianos sobre $X$, mediante las fórmulas:
\begin{equation}\tag{1.1}
  h_A(B)=\operatorname{Hom}_{\Osh}(X;A,B),
  \qquad
  \hsh_A(B)=\Homsh_{\Osh}(A,B).
\end{equation}

Los funtores $h_A$ y $\hsh_A$ son funtores covariantes exactos por la izquierda,
de los cuales consideramos los funtores derivados derechos, denotados respectivamente por
$\operatorname{Ext}_{\Osh}^{p}(X;A,B)$ y $\Extsh_{\Osh}^{p}(A,B)$. Por tanto,
por definición,
\begin{equation}\tag{1.2}
\left\{
\begin{aligned}
\operatorname{Ext}_{\Osh}^{p}(X;A,B)
  &= (R^{p}h_A)(B)
   = H^{p}\bigl(\operatorname{Hom}_{\Osh}(X;A,C(B))\bigr),\\
\Extsh_{\Osh}^{p}(A,B)
  &= (R^{p}\hsh_A)(B)
   = H^{p}\bigl(\Homsh_{\Osh}(A,C(B))\bigr),
\end{aligned}
\right.
\end{equation}

donde el símbolo $R^{p}$ designa el paso a los funtores derivados derechos, y donde
$C(B)$ designa una resolución inyectiva arbitraria de $B$ en $\Ccat^{\Osh}$.
Denotemos por
\[
  \Gamma:\Ccat^{\Zsh}\longrightarrow\Ccat
\]
el funtor \fquote{secciones}. Recordemos que sus funtores derivados derechos se
denotan por
\[
  B\longmapsto H^{p}(X,B).
\]

\src{170}
\begin{equation}\tag{1.3}
  H^{p}(X,B)=(R^{p}\Gamma)(B)=H^{p}\bigl(\Gamma(C(B))\bigr).
\end{equation}
Evidentemente, se tiene $h_A=\Gamma\hsh_A$; por otra parte, se comprueba que $\hsh_A$
transforma objetos inyectivos en objetos acíclicos para $\Gamma$. Se concluye del
modo habitual:

\result{PROPOSICIÓN}{1} Para todo $\Osh$-Módulo $A$ existe un funtor
espectral cohomológico sobre $\Ccat^{\Osh}$, que converge al funtor graduado
$\bigl(\operatorname{Ext}_{\Osh}^{*}(X;A,B)\bigr)$ y cuyo término inicial es
\begin{equation}\tag{1.4}
  E_{2}^{p,q}(A,B)=H^{p}\bigl(X,\Extsh_{\Osh}^{q}(A,B)\bigr).
\end{equation}
De ella se deducen \fquote{homomorfismos de borde} y una sucesión exacta de cinco términos
que no escribiremos.

\result{COROLARIO}{1} Si $A$ es localmente isomorfo a $\Osh^{n}$, entonces se tienen
isomorfismos canónicos
\begin{equation}\tag{1.5}
  \operatorname{Ext}_{\Osh}^{p}(X;A,B)
  \overset{\sim}{\longleftarrow}
  H^{p}\bigl(X,\Homsh_{\Osh}(A,B)\bigr)
\end{equation}
(dados por homomorfismos de borde de la sucesión espectral). En particular, se tiene
un isomorfismo canónico
\begin{equation}\tag{1.6}
  \operatorname{Ext}_{\Osh}^{p}(X;\Osh,B)=H^{p}(X,B).
\end{equation}

Para utilizar estos resultados, es preciso saber explicitar los
$\Extsh_{\Osh}^{p}(A,B)$. Son funtores que se calculan localmente,
es decir, si $U$ es un abierto de $X$, se tiene
\[
  \Extsh_{\Osh}^{p}(A,B)|U
  =\Extsh_{\Osh|U}^{p}(A|U,B|U),
\]
como resulta del hecho de que la restricción a $U$ de un $\Osh$-Módulo inyectivo
es un $(\Osh|U)$-Módulo inyectivo. Además, para $x\in X$ fijo, se tienen homomorfismos funtoriales
\begin{equation}\tag{1.7}
  \Homsh_{\Osh}(A,B)_x
  \longrightarrow
  \operatorname{Hom}_{\Osh_x}(A_x,B_x),
\end{equation}
que se prolongan de manera única a un homomorfismo de funtores cohomológicos
en $B$:
\begin{equation}\tag{1.8}
  \Extsh_{\Osh}^{p}(A,B)_x
  \longrightarrow
  \operatorname{Ext}_{\Osh_x}^{p}(A_x,B_x).
\end{equation}

\result{PROPOSICIÓN}{2} Si $A$ es, en un entorno de $x$, isomorfo al
conúcleo de un homomorfismo $\Osh^{m}\longrightarrow\Osh^{n}$, entonces (1.7) es un
isomorfismo para todo $p$.\footnote{La rama diplomática conserva la lectura
impresa « (1.7) » y la afirmación en todos los grados; la rama actual corrige
la referencia a « (1.8) » y restringe el alcance a las hipótesis que justifican
cada grado.} Esto ocurre, en
particular, si $A$ es un $\Osh$-Módulo coherente [3].

\src{171}
\result{PROPOSICIÓN}{3} Sea $L_{*}=(L_i)$ una resolución por la izquierda del
$\Osh$-Módulo $A$ por $\Osh$-Módulos cada uno de los cuales es localmente isomorfo a
un $\Osh^{n}$. Entonces $\Extsh_{\Osh}^{*}(A,B)$ se identifica con
$H^{*}\bigl(\Homsh_{\Osh}(L_{*},B)\bigr)$, y
$\operatorname{Ext}_{\Osh}^{*}(X;A,B)$ se identifica con la hipercohomología de $X$
con respecto al complejo $\Homsh_{\Osh}(L_{*},B)$.

La demostración es estándar: se considera el bicomplejo
$\Homsh_{\Osh}(L_{*},C(B))$, donde $C(B)$ es una resolución inyectiva de $B$, y los
homomorfismos naturales de $\Homsh_{\Osh}(L_{*},B)$ y de
$\Homsh_{\Osh}(A,C(B))$ en este último.

Para terminar, señalemos que los dos funtores $\operatorname{Ext}$ introducidos
en (1.2) no solo son funtores cohomológicos en $B$, sino también
bifuntores cohomológicos, covariantes en $B$ y contravariantes en $A$.

\fsection{2}{Ley de composición en los Ext}

Los resultados de este número se deben, independientemente, a CARTIER y a YONEDA; véase
una exposición de CARTIER [1] para más detalles. Sea $\Ccat$ una categoría abeliana.
Sean $K$ y $L$ dos objetos graduados de $\Ccat$. Denotamos por
$\operatorname{Hom}(K,L)$ el grupo abeliano graduado cuya componente de grado
$n$ está formada por los homomorfismos homogéneos de grado $n$ de $K$ en $L$ (es decir,
los sistemas $(u_i)$ de homomorfismos $K^{i}\longrightarrow L^{i+n}$).

Si $K$ y $L$ son complejos (con operadores diferenciales de grado $+1$, para
fijar ideas), se introduce en $\operatorname{Hom}(K,L)$ el operador
diferencial dado por
\begin{equation}\tag{2.1}
  \delta(u)=du+(-1)^{n+1}ud,\qquad\text{donde }n=\deg(u),
\end{equation}
que lo convierte en un complejo con operador diferencial de grado $+1$. Los ciclos de
grado $n$ son los operadores de grado $n$ que anticonmutan con $u$ (como
operadores homogéneos).\footnote{La rama diplomática conserva la lectura
impresa « con $u$ »; la rama actual la corrige a « con $d$ » conforme a
la ecuación (2.1).} Se puede entonces considerar
$H^{*}\bigl(\operatorname{Hom}(K,L)\bigr)$; este es un invariante de los tipos
de homotopía de $K$ y de $L$, que podremos denotar por $H^{*}(K,L)$.

Si se tiene un tercer complejo $M$, la composición de los homomorfismos
define un acoplamiento
\[
  \operatorname{Hom}(K,L)\times\operatorname{Hom}(L,M)
  \longrightarrow\operatorname{Hom}(K,M)
\]
compatible con los operadores diferenciales, de donde, al pasar a la cohomología,
se obtienen acoplamientos
\begin{equation}\tag{2.2}
  H^{*}(K,L)\times H^{*}(L,M)\longrightarrow H^{*}(K,M),
\end{equation}
denotados por $(u,v)\longmapsto vu$. Estos acoplamientos satisfacen la evidente propiedad
de asociatividad. En particular, $H^{*}(K,K)$ es un anillo graduado asociativo
con unidad, y $H^{*}(K,L)$ (resp. $H^{*}(L,K)$) es un módulo graduado por la derecha
(resp. por la izquierda) sobre este anillo, etc. En grado $0$, (2.2) se reduce a la
composición de los morfismos de complejos. Finalmente, una sucesión exacta de
complejos
\[
  0\longrightarrow K'\longrightarrow K\longrightarrow K''\longrightarrow 0,
\]
\src{172}
tal que $K'^{i}$ se identifica con un sumando directo de $K^{i}$ para todo $i$,
da lugar a una sucesión exacta de complejos de grupos
$\operatorname{Hom}(K'',L)$, etc., y de ahí a un operador de conexión
\[
  H^{i}(K',L)\longrightarrow H^{i+1}(K'',L).
\]
Se definen de manera análoga los operadores de conexión relativos a una sucesión exacta en $L$. Los
acoplamientos (2.2) son compatibles, en el sentido usual, con estos operadores de conexión.

Supongamos ahora que $\Ccat$ es una categoría tal que todo objeto $A$
de $\Ccat$ admite una resolución inyectiva $C(A)$. Se comprueba entonces, utilizando
una de las numerosas variantes del teorema del bicomplejo, que
\[
  H^{*}\bigl(C(A),C(B)\bigr)
  =H^{*}\bigl(\operatorname{Hom}(C(A),C(B))\bigr)
\]
es canónicamente isomorfo a
\[
  H^{*}\bigl(\operatorname{Hom}(A,C(B))\bigr)
  =\operatorname{Ext}^{*}(A,B).
\]
Los operadores de conexión considerados anteriormente dan los homomorfismos de conexión de los
$\operatorname{Ext}$. Además, los acoplamientos (2.2) dan aquí
acoplamientos asociativos, compatibles con los homomorfismos de conexión:
\begin{equation}\tag{2.3}
  \operatorname{Ext}^{*}(A,B)\times\operatorname{Ext}^{*}(B,C)
  \longrightarrow\operatorname{Ext}^{*}(A,C).
\end{equation}
En particular, $\operatorname{Ext}^{*}(A,A)$ es un anillo graduado asociativo
con unidad, etc. (Se demuestra de manera análoga que los funtores
$\operatorname{Ext}$ operan en los funtores derivados de un funtor
cualquiera; no utilizaremos aquí este hecho.)

En el caso en que la categoría considerada sea la categoría $\Ccat^{\Osh}$ de los
$\Osh$-Módulos sobre $X$, se obtienen, por tanto, acoplamientos
\begin{equation}\tag{2.4}
  \operatorname{Ext}_{\Osh}^{p}(X;A,B)
  \times\operatorname{Ext}_{\Osh}^{q}(X;B,C)
  \longrightarrow\operatorname{Ext}_{\Osh}^{p+q}(X;A,C),
\end{equation}
que pueden explicitarse como se ha indicado. Por lo demás, los mismos
desarrollos, sustituyendo la categoría de los grupos abelianos por la
categoría de los haces abelianos sobre $X$, y los funtores $\operatorname{Hom}$
por los funtores $\Homsh$, definen también acoplamientos con las mismas
propiedades formales y, esta vez, de \fquote{naturaleza local}:
\begin{equation}\tag{2.5}
  \Extsh_{\Osh}^{p}(A,B)\times\Extsh_{\Osh}^{q}(B,C)
  \longrightarrow\Extsh_{\Osh}^{p+q}(A,C).
\end{equation}
Estos últimos se precisan observando que los homomorfismos (1.8) son
compatibles con los acoplamientos entre los $\operatorname{Ext}$.

\src{173}
Recordemos, por último, que existe también una estructura multiplicativa entre los funtores
$H^{p}(X,A)$ (producto cup). Se comprueba entonces que las sucesiones espectrales de la
proposición 1 son compatibles con las estructuras multiplicativas; más
precisamente, se tiene un acoplamiento de la sucesión espectral $E(A,B)$ y de la sucesión
espectral $E(B,C)$ hacia la sucesión espectral $E(A,C)$ que, en el límite,
se reduce al acoplamiento entre los $\operatorname{Ext}$ globales y, en el
término inicial, al acoplamiento deducido, en el segundo miembro de (1.4), del
producto cup y de los acoplamientos de los $\operatorname{Ext}$ locales. De ello resulta,
en particular, que los \fquote{homomorfismos de borde}
\begin{equation}\tag{2.6}
  \operatorname{Ext}_{\Osh}^{n}(X;A,B)
  \longrightarrow H^{0}\bigl(X,\Extsh_{\Osh}^{n}(A,B)\bigr)
\end{equation}
y
\begin{equation}\tag{2.7}
  H^{n}\bigl(X,\Homsh_{\Osh}(A,B)\bigr)
  \longrightarrow\operatorname{Ext}_{\Osh}^{n}(X;A,B)
\end{equation}
son compatibles con las estructuras multiplicativas. Por tanto, si nos restringimos a los
haces localmente isomorfos a un $\Osh^{m}$, esto explicita por completo la
composición entre los $\operatorname{Ext}$ globales mediante el producto cup, teniendo
en cuenta el isomorfismo (1.5).

\fsection{3}{Resultados de cohomología local}

Sea $A$ un anillo conmutativo con unidad provisto de un ideal $J$. Vamos a
definir, para un $A$-módulo variable $M$, homomorfismos funtoriales
\begin{equation}\tag{3.1}
\left\{
\begin{aligned}
\operatorname{Ext}_{A}^{p}(A/J,M)
  &\longrightarrow
  \operatorname{Hom}_{A}\!\left(\bigwedge^{p}J/J^{2},M\otimes A/J\right),\\
\operatorname{Tor}_{p}^{A}(A/J,M)
  &\longleftarrow
  \left(\bigwedge^{p}J/J^{2}\right)
  \otimes\operatorname{Hom}_{A}(A/J,M),
\end{aligned}
\right.
\end{equation}
donde los productos tensoriales y exteriores se toman sobre el anillo $A$; nótese,
por otra parte, que $J/J^{2}$ es en realidad un $A/J$-módulo, y que sus potencias
exteriores como $A$-módulo o como $A/J$-módulo son las mismas.
Definir los homomorfismos (3.1) equivale a definir, para todo
sistema $\underline{x}=(x_1,\ldots,x_p)$ de elementos de $J$, homomorfismos
$\varphi_{\underline{x}}$:
\begin{equation}\tag{3.2}
\left\{
\begin{aligned}
\varphi_{\underline{x}}:
  \operatorname{Ext}_{A}^{p}(A/J,M)&\longrightarrow M\otimes A/J,\\
\varphi_{\underline{x}}:
  \operatorname{Hom}_{A}(A/J,M)&\longrightarrow
  \operatorname{Tor}_{p}^{A}(A/J,M),
\end{aligned}
\right.
\end{equation}
de tal manera que se satisfagan las condiciones siguientes:

\src{174}
\begin{enumerate}
\item[\textup{i.}] $\varphi_{x_1,\ldots,x_p}$ depende de manera
$A$-multilineal alternada del sistema de los $x_i\in J$.
\item[\textup{ii.}] $\varphi_{x_1,\ldots,x_p}$ es nulo cuando alguno de los $x_i$ está
en $J^{2}$.
\end{enumerate}
En realidad, ii. resulta de i., pues $a\varphi_{\underline{x}}=0$ para $a\in J$,
como se ve observando que todos los módulos de (3.2) son anulados por
$J$.

Para definir $\varphi_{\underline{x}}$, consideremos el complejo
$K_{\underline{x}}$ cuyo $A$-módulo subyacente es $\bigwedge A^{p}$, y
cuyo operador diferencial es el producto interior $i_{\underline{x}}$ por
$\underline{x}$ considerado como forma lineal sobre $A^{p}$ de componentes
$x_1,\ldots,x_p$. El operador diferencial es de grado $-1$, los grados son
positivos, y la cohomología de este complejo en grado $0$ es
$A/(x_1A+\cdots+x_pA)$. Como los $x_i$ están en $J$, se deduce una
aumentación $K_{\underline{x},0}\longrightarrow A/J$. Así,
$K_{\underline{x}}$ aparece como un complejo libre aumentado, con módulo
de aumentación $A/J$. De ello se deducen, del modo habitual, homomorfismos
\[
\operatorname{Ext}_{A}^{*}\bigl(H_0(K_{\underline{x}}),M\bigr)
\longrightarrow
H^{*}\bigl(\operatorname{Hom}_{A}(K_{\underline{x}},M)\bigr)
\]
y
\[
\operatorname{Tor}_{*}^{A}\bigl(H_0(K_{\underline{x}}),M\bigr)
\longleftarrow
H_{*}\bigl(K_{\underline{x}}\otimes_A M\bigr),
\]
de donde, componiendo con los homomorfismos sobre los $\operatorname{Ext}$ y los
$\operatorname{Tor}$ deducidos del homomorfismo de aumentación
$H_0(K_{\underline{x}})\longrightarrow A/J$, se obtienen homomorfismos
\begin{equation}\tag{3.3}
\left\{
\begin{aligned}
\psi_{\underline{x}}:
\operatorname{Ext}_{A}^{*}(A/J,M)
&\longrightarrow H^{*}\bigl(\operatorname{Hom}_{A}(K_{\underline{x}},M)\bigr),\\
\psi_{\underline{x}}:
\operatorname{Tor}_{*}^{A}(A/J,M)
&\longleftarrow H_{*}\bigl(K_{\underline{x}}\otimes_A M\bigr).
\end{aligned}
\right.
\end{equation}
Ahora bien, se comprueba inmediatamente que, en grado máximo $p$, la cohomología de los segundos
miembros es $M/(x_1M+\cdots+x_pM)$ (resp. es el conjunto de los elementos de $M$
anulados por cada uno de los $x_i$). Como los $x_i$ están en $J$, se deducen
homomorfismos
\begin{equation}\tag{3.4}
\left\{
\begin{aligned}
H^{p}\bigl(\operatorname{Hom}_{A}(K_{\underline{x}},M)\bigr)
&\longrightarrow M\otimes A/J,\\
H_{p}\bigl(K_{\underline{x}}\otimes_A M\bigr)
&\longleftarrow\operatorname{Hom}_{A}(A/J,M).
\end{aligned}
\right.
\end{equation}
Componiendo los homomorfismos (3.3) y (3.4), se obtienen los homomorfismos (3.2)
que queríamos definir. La comprobación de i. es laboriosa, pero no
presenta dificultades.

\src{175}
\result{PROPOSICIÓN}{4} Sea $A$ un anillo conmutativo con unidad, y sea
$(x_1,\ldots,x_p)$ una sucesión de elementos de $A$ tal que, para $1\leq i\leq p$,
la imagen de $x_i$ en el cociente de $A$ por el ideal generado por
$(x_1,\ldots,x_{i-1})$ no sea divisor de $0$. Sea $J$ el ideal generado
por los $x_i$. Entonces $J/J^{2}$ es un $(A/J)$-módulo libre que tiene por base las
imágenes canónicas de los $x_i$, el complejo $K_{\underline{x}}$ es una resolución
libre de $A/J$, y para todo $A$-módulo $M$, los homomorfismos (3.1) en
grado $p$ son biyectivos. Lo mismo ocurre con los homomorfismos análogos
definidos para grados $i$ cualesquiera, siempre que $J.M=0$.

(El punto esencial, del cual resultan todos los demás, es la aciclicidad de
$K_{\underline{x}}$, que es un hecho bien conocido bajo las condiciones indicadas.)

\result{COROLARIO}{1} Sean $A$ y $J$ como arriba, y supongamos además que
$A$ es un álgebra afín regular de $\dim n$ sobre un cuerpo perfecto $k$, y
que $A/J$ es un álgebra afín regular. Denotemos por $\Omega^{1}(A)$,
$\Omega^{1}(A/J)$ los módulos de diferenciales de Kähler. Entonces se tiene un
isomorfismo canónico compatible con la localización,
\begin{equation}\tag{3.5}
  \operatorname{Ext}_{A}^{p}\bigl(\Omega^{n-p}(A/J),\Omega^{n}(A)\bigr)=A/J.
\end{equation}
En efecto, $\Omega^{n-p}(A/J)$ es un $(A/J)$-módulo libre de rango $1$, y asimismo
$\Omega^{n}(A)$ es un $A$-módulo libre de rango $n$,\footnote{La rama diplomática
conserva las dos lecturas impresas « libre » y la segunda lectura « rango $n$ »;
la rama actual las corrige a « proyectivo de rango $1$ », como exige también
la identificación con $\bigwedge^{p}(J/J^{2})$ más abajo.}
por lo que el primer miembro es igual a
\[
  \operatorname{Ext}_{A}^{p}(A/J,A)
  \otimes\Omega^{n-p}(A/J)'\otimes\Omega^{n}(A)
\]
(donde el símbolo « $'$ » designa el $(A/J)$-módulo dual). El producto tensorial de los
dos últimos factores se identifica con $\bigwedge^{p}(J/J^{2})$, por lo que el conjunto
se identifica con
\[
  \operatorname{Ext}_{A}^{p}\!\left(A/J,\bigwedge^{p}(J/J^{2})\right),
\]
y, por la proposición, con
\[
  \operatorname{Hom}_{A}\!\left(\bigwedge^{p}J/J^{2},
  \bigwedge^{p}J/J^{2}\right),
\]
es decir, con $A/J$. En particular, en
$\operatorname{Ext}_{A}^{p}\bigl(\Omega^{n-p}(A/J),\Omega^{n}(A)\bigr)$ hay un
elemento distinguido, correspondiente al elemento unidad de $A/J$, llamado
\emph{clase fundamental del ideal $J$ en $A$}. (En realidad, puede
definirse bajo condiciones sensiblemente más generales.) El
corolario 1 puede darse en una forma más geométrica y más global:

\result{COROLARIO}{2} Sean $X$ una variedad no singular sobre un cuerpo
algebraicamente cerrado $k$, $Y$ una subvariedad cerrada no singular de $X$,
$\Osh_X$ el haz de anillos locales de $X$, y $\Osh_Y$ el haz de anillos
locales de $Y$ considerado como haz cociente de $\Osh_X$. Sean $n$ la
dimensión de $X$ y $n-p$ la de $Y$.

\src{176}
Sea $\Omega_X$ (resp. $\Omega_Y$) el haz de gérmenes de formas
diferenciales regulares sobre $X$ (resp. $Y$). Entonces se tienen isomorfismos
canónicos:
\begin{equation}\tag{3.6}
  \Extsh_{\Osh_X}^{p}\bigl(\Omega_Y^{n-p},\Omega_X^{n}\bigr)=\Osh_Y,
\end{equation}
o también
\begin{equation}\tag{3.6 bis}
  \Extsh_{\Osh_X}^{p}\bigl(\Osh_Y,\Omega_X^{n}\bigr)=\Omega_Y^{n-p}.
\end{equation}
La fórmula (3.6 bis) puede servir como definición de $\Omega_Y^{n-p}$ cuando $Y$
es una variedad singular. Más precisamente:

\result{PROPOSICIÓN}{5} Sea $Y$ un subconjunto algebraico de dimensión
$q=n-p$ en una variedad algebraica no singular $X$ de dimensión $n$. Sean
$F$ un haz algebraico coherente sobre $X$ cuyo soporte está contenido en $Y$, y $L$
un haz algebraico localmente libre sobre $X$. Entonces los haces
$\Extsh_{\Osh_X}^{i}(F,L)$ son nulos para $i<p$, mientras que para $i=p$ se tiene un
isomorfismo canónico
\begin{equation}\tag{3.7}
  \Extsh_{\Osh_X}^{p}(F,L)
  =\Homsh_{\Osh_X}\!\left(F,
  \Extsh_{\Osh_X}^{p}(\Osh_X/J,L)\right),
\end{equation}
donde $J$ designa un haz de ideales arbitrario sobre $X$ que anula $F$ y cuyo
conjunto de ceros es $Y$. En particular, si $F$ es un haz algebraico
coherente sobre $Y$, se tiene
\begin{equation}\tag{3.7 bis}
  \Extsh_{\Osh_X}^{p}(F,L)
  =\Homsh_{\Osh_Y}\!\left(F,\Extsh_{\Osh_X}^{p}(\Osh_Y,L)\right).
\end{equation}
Por último, siendo todavía $F$ un haz algebraico coherente sobre $Y$, los haces
\[
  E^{i}(F)=\Extsh_{\Osh_X}^{p+i}(F,\Omega_X^{n})
\]
no dependen de la manera de sumergir el espacio algebraico $Y$ en una variedad
no singular $X$.

Como la cuestión es local, se puede suponer que $X$ es afín y $L=\Osh_X$. Esto reduce
el problema a una cuestión de álgebra conmutativa, e incluso de álgebra local: si $A$ es un
anillo local regular y $M$ un $A$-módulo cuyo soporte tiene dimensión
$\leq q=n-p$, hay que demostrar que $\operatorname{Ext}_{A}^{i}(M,A)=0$ para $i<p$, y
\[
  \operatorname{Ext}_{A}^{p}(M,A)
  =\operatorname{Hom}_{A}\!\left(M,\operatorname{Ext}_{A}^{p}(A/J,A)\right),
\]
donde $J$ es un ideal cualquiera de « dimensión » $\leq q$ que anula $M$. Para el
primer punto se procede por inducción sobre $q$: un devissage inmediato reduce
al caso en que $M$ es de la forma $A/J$ y luego, sustituyendo $J$ por un ideal más
pequeño y utilizando la hipótesis de inducción y la sucesión exacta de los
$\operatorname{Ext}$, al caso en que $J$ está generado por un « sistema de parámetros »
como en la proposición 4, donde el resultado es inmediato. El resultado
anterior implica que, si $J$ es un ideal

\src{177}
fijo de « dimensión » $\leq q$, entonces el funtor contravariante
$E(M)=\operatorname{Ext}_{A}^{p}(M,A)$ sobre la categoría de los $A/J$-módulos es
exacto por la izquierda; además, transforma sumas directas en productos directos,
de donde resulta fácilmente que
\[
  E(M)=\operatorname{Hom}_{A}\bigl(M,E(A)\bigr).
\]
\footnote{La rama diplomática conserva la lectura impresa $E(A)$; en la rama
actual, el objeto representante es $A/J$, y la igualdad requerida para (3.7)
es $E(A/J)$.} Por último, la
última afirmación de la proposición 5 es más sutil y resulta de una
caracterización intrínseca de los $E^{i}(F)$ mediante un teorema de dualidad
local que no puede exponerse aquí.

\result{COROLARIO}{} Denotemos por $\omega_Y^{q}$ el haz
$\Extsh_{\Osh_X}^{p}(\Osh_Y,\Omega_X^{n})$. Entonces se tiene un isomorfismo
funtorial para los haces algebraicos coherentes $F$ sobre $Y$:
\begin{equation}\tag{3.8}
  \Extsh_{\Osh_X}^{p}(F,\Omega_X^{n})
  =\Homsh_{\Osh_X}(F,\omega_Y^{q}).
\end{equation}

\fsection{4}{Clase de cohomología asociada a una subvariedad}

En todo lo que sigue, $X$ designa un conjunto algebraico de dimensión $n$, definido
sobre un cuerpo $k$, que supondremos algebraicamente cerrado para simplificar. Salvo
en el n\textsuperscript{o} 6, se supone que $X$ es no singular. Denotamos por
$\Osh_X$ el haz de anillos locales sobre $X$, y por
\[
  \Omega_X^{*}=\bigcup_p\Omega_X^{p}
\]
el haz de gérmenes de formas diferenciales sobre $X$.\footnote{La rama
diplomática conserva el signo de unión impreso; la rama actual lo corrige a
la suma directa que define la estructura graduada de los haces $\Omega_X^p$.}
Si $Y$ es una parte cerrada de $X$, identificamos
los haces algebraicos coherentes sobre $Y$ con haces coherentes sobre $X$
nulos fuera de $Y$; en particular, así ocurre con $\Osh_Y$ y $\Omega_Y$.

\result{LEMA}{1} Sean $F$ un haz algebraico coherente sobre $X$ cuyo
soporte tiene dimensión $\leq n-p$, y $L$ un haz algebraico coherente
localmente libre sobre $X$. Entonces $\operatorname{Ext}_{\Osh_X}^{i}(X;F,L)$ es nulo
para $i<p$, y se tiene un isomorfismo canónico
\begin{equation}\tag{4.1}
  \operatorname{Ext}_{\Osh_X}^{p}(X;F,L)
  =H^{0}\bigl(X,\Extsh_{\Osh_X}^{p}(F,L)\bigr).
\end{equation}
Si $F$ es un haz algebraico coherente sobre una parte cerrada $Y$ de $X$ de
dimensión $\leq n-p$, se tiene un isomorfismo canónico
\begin{equation}\tag{4.1 bis}
  \Extsh_{\Osh_X}^{p}(F,L)
  =\Homsh_{\Osh_X}\!\left(
    F\otimes L'\otimes\Omega_X^{n},\omega_Y^{n-p}
  \right),
\end{equation}
donde $\omega_Y^{n-p}$ es el haz sobre $Y$ definido en el corolario de la
proposición 5 (que se identifica con $\Omega_Y^{n-p}$ si $Y$ es no singular).

La fórmula (4.1) es una consecuencia inmediata de la sucesión espectral de la
proposición 1 y de la proposición 5; en virtud de la fórmula (3.8), se puede
escribir

\src{178}
\begin{align*}
\Extsh_{\Osh_X}^{p}(F,L)
&=L\otimes(\Omega_X^{n})'\otimes
  \Extsh_{\Osh_X}^{p}(F,\Omega_X^{n})\\
&=L\otimes(\Omega_X^{n})'\otimes
  \Homsh_{\Osh_X}(F,\omega_Y^{q})\\
&=\Homsh_{\Osh_X}\!\left(
  F\otimes L'\otimes\Omega_X^{n},\omega_Y^{q}\right),
\end{align*}
donde $q=n-p$, de donde resulta la fórmula (4.1 bis).

Tomando, en particular, $F=\Osh_Y$, $L=\Omega_X^{p}$, se obtiene, teniendo en cuenta
que $\Omega_X^{n}\otimes(\Omega_X^{p})'=\Omega_X^{n-p}$, un isomorfismo
canónico
\begin{equation}\tag{4.2}
  \operatorname{Ext}_{\Osh_X}^{p}(X;\Osh_Y,\Omega_X^{p})
  =\operatorname{Hom}_{\Osh_X}
  (X;\Omega_X^{n-p},\omega_Y^{n-p}).
\end{equation}
Supongamos ahora que $Y$ es no singular, para simplificar, de modo que
$\omega_Y^{n-p}=\Omega_Y^{n-p}$. Se tiene un homomorfismo natural de
$\Omega_X^{n-p}$ en $\Omega_Y^{n-p}$, de donde resulta una sección canónica $s_Y$ del
haz $\Extsh_{\Osh_X}^{p}(\Osh_Y,\Omega_X^{p})$, que llamaremos, si
todas las componentes de $Y$ tienen dimensión $n-p$, \emph{sección
fundamental} del haz $\Extsh_{\Osh_X}^{p}(\Osh_Y,\Omega_X^{p})$. Esta
sección define, en virtud de (4.1), un elemento de
$\operatorname{Ext}_{\Osh_X}^{p}(X;\Osh_Y,\Omega_X^{p})$. Ahora bien, el homomorfismo
natural $\Osh_X\longrightarrow\Osh_Y$ define un homomorfismo
\[
  \operatorname{Ext}_{\Osh_X}^{p}(X;\Osh_Y,\Omega_X^{p})
  \longrightarrow
  \operatorname{Ext}_{\Osh_X}^{p}(X;\Osh_X,\Omega_X^{p})
  =H^{p}(X,\Omega_X^{p}).
\]
Se obtiene así un elemento de $H^{p}(X,\Omega_X^{p})$, denotado por $P_X(Y)$ y llamado
\emph{clase de cohomología de $Y$ en $X$}. Se deduce, por tanto, de la sección
$s_Y$ de $\Extsh_{\Osh_X}^{p}(\Osh_Y,\Omega_X^{p})$ mediante el diagrama
de homomorfismos siguiente:
\begin{equation}\tag{4.3}
\begin{aligned}
H^{p}(X,\Omega_X^{p})
&=\operatorname{Ext}_{\Osh_X}^{p}(X;\Osh_X,\Omega_X^{p})\\
&\longleftarrow\operatorname{Ext}_{\Osh_X}^{p}(X;\Osh_Y,\Omega_X^{p})
\xrightarrow{\ \sim\ }
H^{0}\!\left(X,\Extsh_{\Osh_X}^{p}(\Osh_Y,\Omega_X^{p})\right).
\end{aligned}
\end{equation}

Llamemos \emph{ciclo no singular de dimensión $n-p$} a un elemento del grupo
abeliano libre generado por las subvariedades irreducibles no singulares de
dimensión $n-p$ en $X$. Entonces la función $Y\longmapsto P(Y)$ se prolonga a
un homomorfismo del grupo de ciclos no singulares de dimensión $n-p$ sobre $X$
en el grupo $H^{p}(X,\Omega_X^{p})$.

Sean $Z^{n-p}$ y $Z'^{n-p'}$ ciclos no singulares de dimensiones $n-p$ y
$n-p'$; se dice que se \emph{intersecan transversalmente} si toda componente
de $Z$ interseca transversalmente toda componente de $Z'$. Entonces el ciclo $Z.Z'$ está
definido y es un ciclo no singular de dimensión $n-p-p'$. Dicho esto, se tiene el

\src{179}
\result{TEOREMA}{1} Si $Z^{n-p}$ y $Z'^{n-p'}$ son ciclos no singulares
que se intersecan transversalmente, entonces se tiene
\begin{equation}\tag{4.4}
  P_X(Z.Z')=P_X(Z).P_X(Z'),
\end{equation}
donde el producto del segundo miembro es el producto cup
\[
  H^{p}(X,\Omega_X^{p})\times H^{p'}(X,\Omega_X^{p'})
  \longrightarrow H^{p+p'}(X,\Omega_X^{p+p'}).
\]
(Se supone que $X$ es isomorfo a una parte localmente cerrada de un espacio
proyectivo.)

Esta última hipótesis sirve únicamente para poder concluir que todo
haz algebraico coherente sobre $X$ es cociente de un haz algebraico
coherente localmente libre (SERRE) y, por tanto, admite una resolución por la izquierda mediante
haces localmente libres. Para demostrar el teorema 1, se puede suponer que
$Z$ y $Z'$ son subvariedades irreducibles no singulares $Y$ y $Y'$ que se
intersecan transversalmente. Sea $L_{*}$ una resolución por la izquierda de $\Osh_Y$ mediante
haces localmente libres; entonces, en virtud de la proposición 3, el diagrama
de homomorfismos (4.3) se identifica con el diagrama
\[
R^{p}\Gamma(\Omega_X^{p})
\xleftarrow{\alpha}
\mathbb R^{p}\Gamma\!\left(\Homsh_{\Osh_X}(L_{*},\Omega_X^{p})\right)
\xrightarrow[\sim]{\beta}
\Gamma\!\left(H^{p}\bigl(\Homsh_{\Osh_X}(L_{*},\Omega_X^{p})\bigr)\right),
\]
donde $\Gamma$ es el funtor « grupo de secciones » sobre la categoría de los
haces abelianos sobre $X$, y $\mathbb R^{p}\Gamma$ designa
la hipercohomología de grado $p$ de dicho funtor, y $R^{p}\Gamma$ su
$p$-ésimo funtor derivado. Supondremos, para simplificar, que $L_0=\Osh_X$ y que
la aumentación $L_0\longrightarrow\Osh_Y$ es el homomorfismo natural (lo cual
es lícito); entonces $\alpha$ se deduce del homomorfismo de complejos
$\Osh_X\longrightarrow L_{*}$ (considerando $\Osh_X$ como un complejo concentrado
en grado $0$), teniendo en cuenta que $\mathbb R^{p}\Gamma(K)=R^{p}\Gamma(K_0)$ si $K$
es un complejo de haces concentrado en grado $0$. El homomorfismo $\beta$ es un
« homomorfismo de borde » bien conocido. Consideremos un diagrama análogo, relativo a
una resolución localmente libre $L'_{*}$ de $\Osh_{Y'}$, y consideremos el
diagrama conmutativo de acoplamientos:
\begin{equation}\tag{4.5}
\resizebox{\textwidth}{!}{$
\begin{array}{ccc}
R^{p}\Gamma(\Omega_X^{p})
&\longleftarrow
\mathbb R^{p}\Gamma\!\left(\Homsh_{\Osh_X}(L_{*},\Omega_X^{p})\right)
&\xrightarrow{\sim}
\Gamma H^{p}\!\left(\Homsh_{\Osh_X}(L_{*},\Omega_X^{p})\right)
\\[3pt]
\times&\times&\times\\[3pt]
R^{p'}\Gamma(\Omega_X^{p'})
&\longleftarrow
\mathbb R^{p'}\Gamma\!\left(\Homsh_{\Osh_X}(L'_{*},\Omega_X^{p'})\right)
&\xrightarrow{\sim}
\Gamma H^{p'}\!\left(\Homsh_{\Osh_X}(L'_{*},\Omega_X^{p'})\right)
\\[3pt]
\downarrow&\downarrow&\downarrow\\[3pt]
R^{p+p'}\Gamma(\Omega_X^{p+p'})
&\longleftarrow
\mathbb R^{p+p'}\Gamma\!\left(
\Homsh_{\Osh_X}(L_{*}\otimes L'_{*},\Omega_X^{p+p'})\right)
&\xrightarrow{\sim}
\Gamma H^{p+p'}\!\left(
\Homsh_{\Osh_X}(L_{*}\otimes L'_{*},\Omega_X^{p+p'})\right).
\end{array}
$}
\end{equation}

\src{180}
Los acoplamientos de las dos columnas de la derecha se deducen del acoplamiento de
complejos de haces
\[
\Homsh_{\Osh_X}(L_{*},\Omega_X^{p})
\times\Homsh_{\Osh_X}(L'_{*},\Omega_X^{p'})
\longrightarrow
\Homsh_{\Osh_X}(L_{*}\otimes L'_{*},\Omega_X^{p+p'}),
\]
que se define mediante el producto exterior
$\Omega_X^{p}\times\Omega_X^{p'}\longrightarrow\Omega_X^{p+p'}$.
El acoplamiento de la primera columna es el producto cup (relativo al producto
exterior). Afirmo que la última fila de (4.5) se identifica con el diagrama
de isomorfismos análogo a (4.3), donde $Y$ se sustituye por $Y\cap Y'$ y $p$ por
$p+p'$. Para ello basta demostrar que $L\otimes L'$ es una resolución
(evidentemente localmente libre) de $\Osh_{Y\cap Y'}$. En efecto, se tiene
\[
  H_0(L\otimes L')=\Osh_Y\otimes\Osh_{Y'}=\Osh_{Y\cap Y'}
\]
y
\[
  H_i(L\otimes L')=\operatorname{Tor}_{i}^{\Osh_X}(\Osh_Y,\Osh_{Y'})=0
  \qquad(i>0),
\]
por el hecho de que $Y$ y $Y'$ se intersecan transversalmente. El teorema 1 resulta
ahora de la fórmula:
\begin{equation}\tag{4.6}
  s_Y.s_{Y'}=s_{Y.Y'}
\end{equation}
(donde el producto del primer miembro es el de la última columna de (4.5)).
Esta fórmula (4.6), de naturaleza puramente local, se comprueba sin dificultad
tomando para $L_{*}$ y $L'_{*}$ las resoluciones consideradas en la proposición
4. Se demuestra asimismo (mediante una demostración más sencilla) que
$Z\longmapsto P_X(Z)$ es compatible con el producto cartesiano:
\begin{equation}\tag{4.7}
  P_{X\times X'}(Z\times Z')=P_X(Z)\otimes P_{X'}(Z')
\end{equation}
(fórmula válida si $Z$, resp. $Z'$, es un ciclo no singular sobre la variedad
no singular $X$, resp. $X'$, considerando $Z\times Z'$ como un ciclo no
singular sobre $X\times X'$). De (4.4) y (4.7) se deduce que $P_X(Z)$ es también
compatible con la operación « imagen inversa » por un morfismo de variedades no
singulares $f:X\longrightarrow X'$:
\begin{equation}\tag{4.8}
  P_X\bigl(f^{-1}(Z)\bigr)=f^{*}\bigl(P_{X'}(Z')\bigr),
\end{equation}
fórmula válida si $Z'$ es un ciclo no singular sobre $X'$ tal que $f$ sea
« transversal » a $Z'$, es decir, tal que la gráfica de $f$ sea transversal al ciclo
$X\times Z'$ en $X\times X'$.

\src{181}
\result{COROLARIO}{1} Sean $X$, $X'$ dos variedades no singulares
localmente cerradas en un espacio proyectivo, y supongamos que $X'$ es completa. Sean $U$
un ciclo no singular sobre $X\times X'$, $a$ y $b$ dos puntos de $X'$ tales que
$U$ interseca transversalmente los ciclos $X\times(a)$ y $X\times(b)$, y sean $Z$
y $Z'$ los ciclos no singulares sobre $X$ tales que
\[
  Z\times(a)=(X\times(a)).U,
  \qquad
  Z'\times(b)=(X\times(b)).u.
\]
Entonces se tiene
\[
  P_X(Z)=P_X(Z').
\]

En efecto, sea $f_a:X\longrightarrow X\times X'$ definido por
$f_a(x)=(x,a)$; entonces, en virtud de (4.8), se tiene la fórmula
$P(Z)=f_a^{*}(P(U))$, y análogamente $P(Z')=f_b^{*}(P(U))$. Ahora bien, utilizando la fórmula de
Künneth
\[
  H^{*}(X\times X',\Omega_{X\times X'}^{*})
  =H^{*}(X,\Omega_X^{*})\otimes H^{*}(X',\Omega_{X'}^{*})
\]
y el hecho de que $H^{0}(X',\Osh_{X'})$ se reduce a los escalares, se obtiene
fácilmente $f_a^{*}=f_b^{*}$, de donde se sigue la conclusión.

Para todo $x\in X$, $(x)$ es una subvariedad no singular de $X$ de
codimensión $n$, y define por tanto un elemento $\varepsilon_x$ de
$H^{n}(X,\Omega_X^{n})$. Si $X$ es una variedad proyectiva no singular, del
corolario 1 resulta que $\varepsilon_x$ no depende del punto $x$ elegido; se
denota por $\varepsilon_X$ y se llama \emph{clase fundamental} de
$H^{n}(X,\Omega_X^{n})$.

\result{OBSERVACIÓN}{} Para disponer de una teoría satisfactoria, habría que definir
$P_X(Z)$ para ciclos $Z$ cualesquiera y demostrar el teorema 1 para una
intersección propia de ciclos. (En el momento de escribir esta exposición, esto aún no se
ha hecho con toda generalidad). Admitiendo que se haya hecho, el corolario 1
pasa a ser: si $Z$ y $Z'$ son dos ciclos algebraicamente equivalentes, entonces
$P_X(Z)=P_X(Z')$ (afirmación que no parece resultar de lo anterior, incluso si
$Z$ y $Z'$ son no singulares).

\fsection{5}{El teorema de dualidad}

En este número, $X$ designa una variedad proyectiva no singular de dimensión
$n$.

\result{TEOREMA}{2} La clase fundamental $\varepsilon_X$ de
$H^{n}(X,\Omega_X^{n})$ es una base de este espacio vectorial.

(La demostración se dará más adelante). En virtud del teorema anterior, se puede
identificar $H^{n}(X,\Omega_X^{n})$ con el cuerpo $k$. Consideremos ahora los
acoplamientos contemplados en el n\textsuperscript{o} 2, que dan en particular un
acoplamiento

\src{182}
\[
\operatorname{Ext}_{\Osh_X}^{p}(X;\Osh_X,F)
\times\operatorname{Ext}_{\Osh_X}^{n-p}(X;F,\Omega_X^{n})
\longrightarrow
\operatorname{Ext}_{\Osh_X}^{n}(X;\Osh_X,\Omega_X^{n}),
\]
es decir,
\begin{equation}\tag{5.1}
  H^{p}(X,F)
  \times\operatorname{Ext}_{\Osh_X}^{n-p}(X;F,\Omega_X^{n})
  \longrightarrow H^{n}(X,\Omega_X^{n})
\end{equation}
Teniendo en cuenta el teorema 2, este acoplamiento define un homomorfismo
\begin{equation}\tag{5.2}
  \operatorname{Ext}_{\Osh_X}^{n-p}(X;F,\Omega_X^{n})
  \longrightarrow\bigl(H^{p}(X,F)\bigr)'
\end{equation}
Este homomorfismo es funtorial en $F$ y conmuta con los homomorfismos de conexión
relativos a las sucesiones exactas $0\longrightarrow F'\longrightarrow F
\longrightarrow F''\longrightarrow0$.

\result{TEOREMA}{3} El homomorfismo (5.2) es un isomorfismo.

Se recupera, en particular, el resultado de Serre:

\result{COROLARIO}{} Sean $E$ un fibrado vectorial algebraico sobre $X$ y
$\Osh_X(E)$ el haz de gérmenes de secciones regulares de $E$; entonces se tienen
isomorfismos canónicos:
\begin{equation}\tag{5.3}
  \bigl(H^{p}(X,\Osh_X(E))\bigr)'
  =H^{n-p}\bigl(X,\Omega_X^{n}\otimes\Osh_X(E')\bigr)
\end{equation}
Basta aplicar el teorema 3 y el corolario 1 a la proposición 1.

Los teoremas 2 y 3 resultarán del enunciado siguiente:

\noindent\textbf{(D) El homomorfismo siguiente}
\begin{equation}\tag{5.2 bis}
  \operatorname{Ext}_{\Osh_X}^{n-p}(X;F,\Omega_X^{n})
  \longrightarrow\bigl(H^{p}(X,F)\bigr)'\otimes L,
  \qquad L=H^{n}(X,\Omega_X^{n}),
\end{equation}
deducido del acoplamiento (5.1) es un isomorfismo.

Demostremos, en efecto, que (D) implica el teorema 2. Sea
$k_x=\Osh_{(x)}$ el haz de anillos locales de la variedad reducida al punto
$x\in X$; consideremos el homomorfismo canónico
$\Osh_X\longrightarrow k_x$ y el homomorfismo asociado
\begin{equation}\tag{5.3}
  H^{0}(X,\Osh_X)\longrightarrow H^{0}(X,k_x),
\end{equation}
su transpuesto se identifica con el homomorfismo
\begin{equation}\tag{5.4}
  \operatorname{Ext}_{\Osh_X}^{n}(X;k_x,\Omega_X^{n})\otimes L'
  \longrightarrow
  \operatorname{Ext}_{\Osh_X}^{n}(X;\Osh_X,\Omega_X^{n})\otimes L'
\end{equation}
deducido del homomorfismo sobre los $\operatorname{Ext}^{n}$ asociado a
$\Osh_X\longrightarrow k_x$:
\begin{equation}\tag{5.5}
  \operatorname{Ext}_{\Osh_X}^{n}(X;k_x,\Omega_X^{n})
  \longrightarrow
  \operatorname{Ext}_{\Osh_X}^{n}(X;\Osh_X,\Omega_X^{n})
  =H^{n}(X,\Omega_X^{n})
\end{equation}

\src{183}
Como (5.3) es un isomorfismo, lo mismo ocurre
con (5.4), y por tanto también con (5.5). Como $s_{(x)}$ es una base de
$\operatorname{Ext}_{\Osh_X}^{n}(X;k_x,\Omega_X^{n})$ en virtud de (4.2), se
sigue que su imagen $\varepsilon_X$ es una base de
$H^{n}(X,\Omega_X^{n})$.

Queda por demostrar el enunciado (D), que resultará de manera puramente formal de los hechos
elementales resumidos en los lemas siguientes. En ellos se supone que $X$ es una
parte cerrada (singular o no) del espacio proyectivo $P$ de dimensión $r$.
Se utiliza la notación $\Osh_P(m)$ para el haz sobre $P$ denotado por $\Osh(m)$ en
[3], y la notación $\Osh_X(m)$ para el haz análogo sobre $X$.

\result{LEMA}{2} El enunciado (D) es verdadero si $X=P$ y $F=\Osh_P(m)$.

Este lema se comprueba mediante un cálculo directo. El cálculo explícito de los
$H^{i}(P,\Osh_P(m))$ se encuentra en [3], pero puede hacerse de manera más
elemental. El cálculo del producto cup
\[
  H^{i}(P,\Osh_P(m))\times H^{j}(P,\Osh_P(m'))
  \longrightarrow H^{i+j}(P,\Osh_P(r+r'))
\]
necesario para explicitar el acoplamiento (5.1) no presenta dificultades.

\result{LEMA}{3} Todo haz algebraico coherente $F$ sobre $X$ es isomorfo a
un haz cociente de un haz $\Osh_X(-m)^{k}$, donde se puede suponer que $m$ es
tan grande como se quiera.

Resulta del hecho de que $F\otimes\Osh_X(m)$ está « generado por sus secciones » para
$m$ grande, cf. [3].

\result{LEMA}{4} Sea $i>0$. Entonces
$H^{r-i}(P,\Osh_P(-m))=0$ para $m$ grande; y para todo haz algebraico
coherente $B$ sobre $X$, se tiene
\[
  \operatorname{Ext}_{\Osh_X}^{i}(X;\Osh_X(-m),B)=0
\]
para $m$ grande.

El primer enunciado resulta de los cálculos explícitos mencionados anteriormente; para el
segundo, se observa que se tiene un isomorfismo
\[
  \operatorname{Ext}_{\Osh_X}^{i}(X;\Osh_X(-m),B)
  =H^{i}(X,B\otimes\Osh_X(m))
\]
(proposición 1, corolario 1), de donde se sigue la conclusión en virtud de un resultado bien
conocido de [3]. Combinando los dos lemas anteriores, se obtiene el

\result{COROLARIO}{} Sea $i>0$. Entonces el funtor
$F\longmapsto H^{r-i}(P,F)$ sobre la categoría de los haces algebraicos coherentes
sobre $P$ es borrable, y lo mismo ocurre con el funtor
$\operatorname{Ext}_{\Osh_X}^{i}(X;F,B)$ sobre la categoría de los haces
algebraicos coherentes sobre $X$.

\src{184}
\result{LEMA}{5} Sean $A$ y $B$ dos haces algebraicos coherentes sobre
$X$, y pongamos $A(m)=A\otimes\Osh_X(m)$. Entonces, para $m$ suficientemente grande,
el homomorfismo canónico
\[
  \operatorname{Ext}_{\Osh_X}^{1}(X;A(-m),B)
  \longrightarrow H^{0}\bigl(X,\Extsh_{\Osh_X}^{1}(A(-m),B)\bigr)
  =H^{0}\bigl(X,\Extsh_{\Osh_X}^{1}(A,B)(m)\bigr)
\]
es un isomorfismo.

Esto resulta inmediatamente de la sucesión espectral de la proposición 1 aplicada a
$A(-m)$ y $B$, pues se tendrá
\[
  E_{2}^{p,q}(A(-m),B)
  =H^{p}\bigl(X,\Extsh_{\Osh_X}^{q}(A(-m),B)\bigr)
  =H^{p}\bigl(X,\Extsh_{\Osh_X}^{q}(A,B)(m)\bigr),
\]
que es nulo para $p>0$ y $m$ grande.

Demostremos ahora (D) en el caso en que $X=P$. Probamos primero que (5.2 bis)
es un isomorfismo para $p=n$; como los dos miembros son entonces
funtores exactos por la izquierda (pues $H^{r+1}(P,F)=0$), del lema 3 resulta
que basta demostrar la afirmación para $F=\Osh_P(-m)$, pero esta se encuentra
contenida en el lema 2. Como los homomorfismos (5.2 bis) son funtoriales y
compatibles con los operadores de conexión, y como para $p<n$ los dos miembros de
(5.2 bis) son funtores borrables en $F$ (corolario del lema 4), se
sigue entonces, mediante un razonamiento estándar, que (5.2 bis) es un isomorfismo
para todo $p$. Esto demuestra el teorema de dualidad para el espacio proyectivo.
Supongamos ahora que $X$ es cualquiera, pero no singular. Por el teorema de
dualidad para $P$, se tiene un isomorfismo
\[
  H^{n}(X,F)'=H^{n}(P,F)'
  =\operatorname{Ext}_{\Osh_P}^{r-n}(P;F,\Omega_P^{r}).
\]
En virtud del lema 1 (n\textsuperscript{o} 4), el último miembro se identifica con
\[
  \operatorname{Hom}_{\Osh_P}(P;F,\omega_X^{n})
  =\operatorname{Hom}_{\Osh_X}(X;F,\Omega_X^{n})
  =\operatorname{Ext}_{\Osh_X}^{0}(X;F,\Omega_X^{n}),
\]
de donde se obtiene un isomorfismo
\begin{equation}\tag{5.6}
  H^{n}(X,F)'=\operatorname{Hom}_{\Osh_X}(X;F,\Omega_X^{n})
  =\operatorname{Ext}_{\Osh_X}^{0}(X;F,\Omega_X^{n}).
\end{equation}
Tomando $F=\Omega_X^{n}$, se obtiene un isomorfismo
\begin{equation}\tag{5.7}
  \gamma:H^{n}(X,\Omega_X^{n})\xrightarrow{\sim}k.
\end{equation}

\src{185}
Se comprueba que el isomorfismo (5.6) no es otro que (5.2 bis) para $p=n$, cuando
se toma $L=k$ gracias a (5.7). Por consiguiente, (5.2 bis) es un isomorfismo para
$p=n$. Para demostrar que es un isomorfismo para todo $p$, basta nuevamente
demostrar que, para $p<n$, los dos miembros de (5.2 bis) son funtores
borrables en $F$, y a fortiori (teniendo en cuenta el lema 3) que ambos miembros
son nulos cuando se toma $F=\Osh_X(-m)$ con $m$ grande. Ahora bien, para el primer
miembro, esto es cierto en virtud del lema 4, y para el segundo miembro se escribe,
utilizando el teorema de dualidad sobre $P$:
\[
  H^{p}(X,\Osh_X(-m))'
  =\operatorname{Ext}_{\Osh_P}^{r-p}
    (P;\Osh_X(-m),\Omega_P^{r}).
\]
El segundo miembro es nulo para $p<n$ y $m$ grande, como resulta del lema 5
(donde se toma $X=P$) y del hecho de que $\Osh_X$, como haz algebraico
coherente sobre $P$, tiene dimensión cohomológica $\leq r-n$ (pues $X$ es no
singular), de donde
\[
  \Extsh_{\Osh_P}^{r-p}(\Osh_X,\Omega_P^{r})=0
  \qquad\text{para }p<n.
\]

\fsection{6}{El teorema de dualidad para las variedades singulares}

Sea $X$ una parte cerrada de dimensión $n$ del espacio proyectivo $P$ de
dimensión $r$. La fórmula (5.6) debe escribirse aquí
\begin{equation}\tag{6.1}
  H^{n}(X,F)'
  \simeq\operatorname{Hom}_{\Osh_X}(X;F,\omega_X^{n})
  =\operatorname{Ext}_{\Osh_X}^{0}(X;F,\omega_X^{n}),
\end{equation}
donde se ha puesto
\begin{equation}\tag{6.2}
  \omega_X^{n}=E^{0}(\Osh_X)
  =\Extsh_{\Osh_P}^{r-n}(\Osh_X,\Omega_P^{r}).
\end{equation}
Como se dijo en la proposición 5, el haz así introducido no depende
en realidad de la inmersión elegida de $X$ en una variedad no singular $P$.
Tomando, en (6.1), $F=\omega_X^{n}$, se obtiene
\begin{equation}\tag{6.2}
  H^{n}(X,\omega_X^{n})'
  \simeq\operatorname{Hom}_{\Osh_X}(X;\omega_X^{n},\omega_X^{n}),
\end{equation}
\footnote{La rama diplomática conserva la segunda numeración impresa « (6.2) »;
la rama actual usa « (6.2 bis) » para distinguir las dos apariciones.} de donde
resulta la existencia de un elemento
distinguido en $H^{n}(X,\omega_X^{n})'$, correspondiente al homomorfismo
identidad de $\omega_X^{n}$ en sí mismo:
\begin{equation}\tag{6.3}
  \gamma:H^{n}(X,\omega_X^{n})\longrightarrow k.
\end{equation}

\src{186}
Consideremos entonces los acoplamientos definidos por la composición de los $\operatorname{Ext}$:
\begin{equation}\tag{6.4}
  H^{p}(X,F)\times
  \operatorname{Ext}_{\Osh_X}^{n-p}(X;F,\omega_X^{n})
  \longrightarrow H^{n}(X,\omega_X^{n}),
\end{equation}
y compongámoslos con el homomorfismo $\gamma$ de (6.3); se obtienen
homomorfismos funtoriales, compatibles con los operadores de borde:
\begin{equation}\tag{6.5}
  \operatorname{Ext}_{\Osh_X}^{n-p}(X;F,\omega_X^{n})
  \longrightarrow H^{p}(X,F)'
\end{equation}
(que generalizan (5.2)). Se comprueba que para $p=n$ se obtiene así
el isomorfismo (6.1). Sentado esto, se tiene el

\result{TEOREMA}{3 bis} Las cuatro condiciones siguientes sobre $X$ son
equivalentes para un entero dado $k\geq0$:
\begin{enumerate}
  \item[\textit{i.}] El homomorfismo funtorial (6.5) es un isomorfismo para
    $n-k\leq p\leq n$.
  \item[\textit{ii.}] Se tiene $H^{p}(X,\Osh_X(-m))=0$ para $m$ grande y para
    $n-k\leq p<n$.
  \item[\textit{iii.}] El funtor $H^{p}(X,F)$ sobre la categoría de los haces
    algebraicos coherentes sobre $X$ es borrable para $n-k\leq p<n$.
  \item[\textit{iv.}] Se tiene
    $E^{i}(\Osh_X)=\Extsh_{\Osh_P}^{r-n+i}(\Osh_X,\Omega_P^{r})=0$
    para $0<i\leq k$.
\end{enumerate}

\noindent\textbf{DEMOSTRACIÓN.} --- \textit{i.}$\Rightarrow$\textit{ii.} en
virtud del lema 4, \textit{ii.}$\Rightarrow$\textit{iii.} en virtud del lema 3,
\textit{iii.}$\Rightarrow$\textit{i.} mediante un razonamiento estándar bien conocido,
teniendo en cuenta que los dos miembros de (6.5) son entonces funtores borrables
para $n-k\leq p<n$ (el primero lo es en virtud del lema 4). Por último, demostremos
\textit{ii.}$\Leftrightarrow$\textit{iv.} Esto resulta del corolario de la
proposición siguiente:

\result{PROPOSICIÓN}{6} Sea $F$ un haz algebraico coherente sobre $X$, y
sea $i$ un entero. Entonces, para $m$ suficientemente grande, se tiene un isomorfismo:
\begin{equation}\tag{6.6}
  H^{i}(X,F(-m))'
  \simeq H^{0}\bigl(X,E^{n-i}(F)(m)\bigr),
\end{equation}
donde se pone (compárese n\textsuperscript{o} 3, proposición 5):
\begin{equation}\tag{6.7}
  E^{j}(F)=\Extsh_{\Osh_P}^{r-n+j}(F,\Omega_P^{r}).
\end{equation}
En efecto, en virtud del teorema de dualidad sobre $P$, el primer miembro de (6.6)
es isomorfo a
$\operatorname{Ext}_{\Osh_P}^{r-i}(P;F(-m),\Omega_P^{r})$, por lo que (6.6) resulta
del lema 5.

\src{187}
\result{COROLARIO}{} Para que se tenga $H^{i}(X,F(-m))=0$ para $m$ grande, es
necesario y suficiente que $E^{n-i}(F)=0$.

Recordemos que los $E^{j}(F)$ no dependen en realidad de la inmersión proyectiva
considerada. La condición del corolario es de naturaleza puramente local; por consiguiente,
si se verifica para $F$, se verifica para todo haz localmente
isomorfo a un haz $F^{n}$. En particular, si esta condición se verifica
para $\Osh_X$, se verifica para todo haz algebraico coherente localmente
libre. Este es, por ejemplo, el caso para todo $i<n$ si $X$ es no singular, para
$i=0$ si ninguna componente de $X$ se reduce a un solo punto, y para $i=0,1$
si $S$ es normal y todas sus componentes tienen dimensión $>1$ (cf.
[3]).\footnote{La rama diplomática conserva la lectura impresa « $S$ es
normal »; la rama actual sustituye la letra porque el párrafo solo se refiere
a la variedad $X$.} Para que esto ocurra para todo
$i<k$, es necesario y suficiente, por definición, que los anillos locales
$\Osh_x$ ($x\in X$) tengan « codimensión homológica » $\geq k$ (cf. [4]
para esta noción). Si $k=n$, esto significa, en virtud del teorema 3 bis, que el
teorema de dualidad es verdadero para $X$, es decir, que (6.5) es un isomorfismo para
todo $p$ y todo $F$. Se pueden dar numerosas condiciones equivalentes sobre
los anillos locales $\Osh_x$ para que esto ocurra (Nagata), por ejemplo que
satisfagan el teorema de equidimensionalidad de Cohen--Macaulay. Así ocurre,
por ejemplo, si localmente $X$ es una « intersección completa » en una
variedad ambiente no singular.

\fsection{7}{La dualidad de Poincaré}

Sea $X$ una variedad proyectiva no singular de dimensión $n$. Entonces
\[
  H^{**}(X)=H^{*}(X,\Omega_X^{*})
\]
es un álgebra anticonmutativa bigraduada de dimensión finita, que graduamos
por el grado total, de modo que $H^{p,q}(X)=H^{p}(X,\Omega_X^{q})$ tiene grado
$p+q$. Los grados de $H^{*}(X)$ están comprendidos entre $0$ y $2n$. En virtud del
teorema 2 y del corolario del teorema 3, $H^{*}(X)$ es un « álgebra de
Poincaré » de dimensión $2n$, es decir, $H^{2n}(X)$ está provisto de un isomorfismo con
el cuerpo base $k$, y el producto
\[
  H^{m}(X)\times H^{2n-m}(X)\longrightarrow H^{2n}(X)=k
\]
es una dualidad entre $H^{m}(X)$ y $H^{2n-m}(X)$. Además, si $Y$ es otra
variedad proyectiva no singular, la fórmula de Künneth para los
haces algebraicos coherentes da
\begin{equation}\tag{7.1}
  H^{*}(X\times Y)=H^{*}(X)\otimes H^{*}(Y),
\end{equation}
isomorfismo compatible con las estructuras de álgebras de Poincaré. Además,
$H^{*}(X)$ es, como álgebra conmutativa, un funtor contravariante en
$X$, pues un morfismo $f:Y\longrightarrow X$ define de manera evidente un
homomorfismo de álgebras graduadas

\src{188}
\begin{equation}\tag{7.2}
  f^{*}:H^{*}(X)\longrightarrow H^{*}(Y).
\end{equation}
Como se trata de álgebras de Poincaré, se obtiene por transposición un
homomorfismo de espacios vectoriales:
\begin{equation}\tag{7.3}
  f_{*}:H^{*}(Y)\longrightarrow H^{*}(X).
\end{equation}
Se vio en el n\textsuperscript{o} 4 que el efecto de $f^{*}$ sobre las clases de
cohomología correspondientes a ciclos no singulares se interpreta
geométricamente tomando las clases de cohomología que corresponden a sus
imágenes inversas. En el caso actual importa demostrar que (7.3) corresponde
análogamente a la operación « imagen directa » de ciclos. Esto resulta, al menos bajo
condiciones adecuadas de no singularidad, del caso particular siguiente:

\result{TEOREMA}{4} Si $f$ es la aplicación identidad de una subvariedad no
singular $Y^{m}$ de $X^{n}$ en $X^{n}$, entonces, denotando por $1_Y$
el elemento unidad de $H(Y)$, se tiene
\begin{equation}\tag{7.4}
  f_{*}(1_Y)=P_X(Y),
\end{equation}
donde el segundo miembro es la clase de cohomología en $X$ asociada a $Y$.

Esta fórmula equivale a
\begin{equation}\tag{7.4 bis}
  \bigl\langle\xi^{m,m},P_X(Y)\bigr\rangle\varepsilon_Y
  =f_{*}(\xi^{m,m})
  \qquad
  \bigl(\xi^{m,m}\in H^{m}(X,\Omega_X^{m})\bigr),
\end{equation}
\footnote{La rama diplomática conserva la fórmula impresa
$f_{*}(\xi^{m,m})$; la rama actual la corrige a $f^{*}(\xi^{m,m})$, pues
$\xi^{m,m}\in H^{m}(X,\Omega_X^{m})$ y ambos miembros pertenecen a
$H^{m}(Y,\Omega_Y^{m})$.} donde
$\varepsilon_Y$ es el elemento fundamental en $H^{m}(Y,\Omega_Y^{m})$, y
constituye, en el caso de las variedades proyectivas no singulares, una nueva
definición de la clase de cohomología asociada a $Y$. Para demostrar el
teorema 4, se interpreta, gracias al teorema 3, el transpuesto del homomorfismo
\[
  H^{m}(X,\Omega_X^{m})
  \longrightarrow H^{m}(Y,\Omega_Y^{m})
  =H^{m}(X,\Omega_Y^{m})
\]
como el homomorfismo
\begin{equation}\tag{7.5}
\begin{split}
  H^{n-m}(X,\Omega_X^{n-m})
  &\simeq
  \operatorname{Ext}_{\Osh_X}^{n-m}(X;\Omega_X^{m},\Omega_X^{n})\\
  &\longleftarrow
  \operatorname{Ext}_{\Osh_X}^{n-m}(X;\Omega_Y^{m},\Omega_X^{n})
  \simeq
  \operatorname{Hom}_{\Osh_X}(X;\Omega_Y^{m},\Omega_Y^{m}).
\end{split}
\end{equation}
Se comprueba que el elemento $1_Y$ del dual de $H^{m}(Y,\Omega_Y^{m})$ se identifica
con el elemento del miembro de la derecha correspondiente al endomorfismo identidad de
$\Omega_Y^{m}$, y por otra parte que la imagen de este elemento en
$H^{n-m}(X,\Omega_X^{n-m})$ es precisamente $P_X(Y)$. Estas compatibilidades, que podrían
haberse dado ya en el n\textsuperscript{o} 4, pueden enunciarse, y son
en realidad válidas, para variedades no singulares cualesquiera; la segunda

\src{189}
resulta, por ejemplo, de la conmutatividad del diagrama siguiente de endomorfismos
canónicos:
\begin{equation}\tag{7.6}
\resizebox{\textwidth}{!}{$
\left\{
\begin{array}{ccccc}
\operatorname{Ext}_{\Osh_X}^{n-m}(X;\Omega_X^{m},\Omega_X^{n})
&\longleftarrow&
\operatorname{Ext}_{\Osh_X}^{n-m}(X;\Omega_Y^{m},\Omega_X^{n})
&\simeq&
\operatorname{Hom}_{\Osh_X}(X;\Omega_Y^{m},\Omega_Y^{m})\\[3pt]
\Big\downarrow&&&&\Big\downarrow\\[3pt]
\operatorname{Ext}_{\Osh_X}^{n-m}(X;\Osh_X,\Omega_X^{n-m})
&\longleftarrow&
\operatorname{Ext}_{\Osh_X}^{n-m}(X;\Osh_Y,\Omega_X^{n-m})
&\simeq&
\operatorname{Hom}_{\Osh_X}(X;\Omega_X^{m},\Omega_Y^{m})
\end{array}
\right.
$}
\end{equation}

Se obtiene así un equivalente exacto del formalismo de la dualidad de Poincaré
sobre las variedades orientadas compactas. En particular, el teorema 4 permite
determinar la clase de cohomología asociada a la diagonal de $X\times X$. Mediante
un razonamiento bien conocido, se deduce, por ejemplo, una fórmula de Lefschetz:

\result{TEOREMA}{5} Sea $f$ un endomorfismo de una variedad proyectiva no
singular $X$, tal que los puntos fijos de $f$ tengan multiplicidad 1.
Entonces el número de estos puntos fijos es congruente, módulo la característica de
$k$, con la suma alternada de las trazas de los endomorfismos de los $H^{i}(X)$ definidos
por $f$.

La restricción sobre $f$ que ha sido necesario imponer se debe a las dificultades mencionadas en
la observación del n\textsuperscript{o} 4. No obstante, se observará que la fórmula de
Lefschetz sigue siendo verdadera si $f$ es « homotópico » a un endomorfismo cuyos
puntos fijos tienen todos multiplicidad 1.

\fsection{8}{Generalización del teorema de dualidad}

Sea $X$ una variedad algebraica no singular tal que todo haz
algebraico coherente $F$ sobre $X$ sea isomorfo a un cociente de un haz
algebraico coherente localmente libre (lo cual ocurre si $X$ es localmente
cerrado en un espacio proyectivo). Entonces todo haz algebraico coherente $F$ sobre
$X$ admite una resolución finita $L$ mediante haces localmente libres, y para
dos resoluciones de esta clase siempre se puede encontrar una tercera y un
homomorfismo de esta en las dos primeras compatible con las
aumentaciones. Análogamente, si $L$ es una resolución finita localmente libre de
$F$, y si se da un homomorfismo $F'\longrightarrow F$, entonces existe
una resolución finita localmente libre $L'$ de $F'$ y un homomorfismo
$L'\longrightarrow L$ compatible con $F'\longrightarrow F$, que se puede
suponer sobreyectivo si $F'\longrightarrow F$ lo es. Esto permite definir, dados
dos enteros $r,s\geq0$, dos multifuntores cohomológicos en los
argumentos $A_1,\ldots,A_r$; $B_1,\ldots,B_s$,

\src{190}
que varían en la categoría de los haces algebraicos coherentes sobre $X$, y cuyos
valores están en la categoría de los haces algebraicos coherentes sobre $X$,
resp. en la categoría de los módulos sobre $H^{0}(X,\Osh_X)$, mediante las fórmulas:
\begin{equation}\tag{8.1}
\resizebox{\textwidth}{!}{$
\left\{
\begin{aligned}
\underline{T}_{r}^{s*}(A_1,\ldots,A_r;B_1,\ldots,B_s)
  &={}H^{*}\!\left(
    \Homsh_{\Osh_X}\!\left(
      L(A_1)\otimes\cdots\otimes L(A_r),
      L(B_1)\otimes\cdots\otimes L(B_s)
    \right)\right),\\[3pt]
T_{r}^{s*}(A_1,\ldots,A_r;B_1,\ldots,B_s)
  &={}\mathbf R^{*}\Gamma\!\left(
    \Homsh_{\Osh_X}\!\left(
      L(A_1)\otimes\cdots\otimes L(A_r),
      L(B_1)\otimes\cdots\otimes L(B_s)
    \right)\right).
\end{aligned}
\right.
$}
\end{equation}
En esta fórmula, $L(F)$ designa una resolución finita localmente libre del
haz algebraico coherente $F$, y $\mathbf R^{*}\Gamma(K)$ designa
la hipercohomología del espacio $X$ con respecto al complejo de haces $K$. Si
$r$ o $s$ es nulo, se sustituye el producto tensorial de los $L(A_i)$, resp. de los
$L(B_j)$, por $\Osh_X$. En particular, $\underline{T}_{0}^{0}$ y $T_{0}^{0}$
son funtores graduados de 0 argumentos, $\underline{T}_{0}^{0}$ se reduce al
grado 0, donde es el haz $\Osh_X$, mientras que $T_{0}^{0}$ es idéntico a
$H^{*}(X,\Osh_X)$. El hecho de que los segundos miembros de (8.1) no dependan
de las resoluciones elegidas es, en cualquier caso, evidente para la primera línea
(pues la cuestión es entonces local), y para la segunda resulta de las observaciones
generales precedentes, teniendo en cuenta la sucesión espectral de
hipercohomología del complejo de haces
\[
  K=\Homsh_{\Osh_X}\bigl(L(A_1)\otimes\cdots,
      L(B_1)\otimes\cdots\bigr),
\]
que converge a la hipercohomología de $X$ con respecto a $K$ y cuyo término
inicial es $H^{p}(X,H^{q}(K))$, es decir,
\begin{equation}\tag{8.2}
  E_{2}^{p,q}
  =H^{p}\!\left(X,
    (\underline{T}_{r}^{s})^{(q)}
      (A_1,\ldots,A_r;B_1,\ldots,B_s)\right).
\end{equation}
Se ve entonces que esta sucesión espectral tampoco depende de las resoluciones
elegidas; su límite es el primer miembro de (8.1). Se definen
fácilmente los operadores de conexión relativos a los diversos argumentos $A_i$, $B_j$,
observando que toda sucesión exacta
$0\longrightarrow F'\longrightarrow F\longrightarrow F''\longrightarrow0$
puede resolverse mediante una sucesión exacta de complejos finitos localmente libres.

En el sistema de los funtores $\underline{T}_{r}^{s*}$, resp.
$T_{r}^{s*}$, se definen operaciones análogas a las del cálculo tensorial, cuya
definición es inmediata a partir de las fórmulas definitorias (8.1). Así, se tiene
una composición (que generaliza la considerada en el n\textsuperscript{o} 2):
\begin{equation}\tag{8.3}
\resizebox{\textwidth}{!}{$
\begin{aligned}
&T_{r}^{s*}(A_1,\ldots,A_r;B_1,\ldots,B_s)
\times T_{r'}^{s'*}(A'_1,\ldots,A'_{r'};B'_1,\ldots,B'_{s'})\\
&\qquad\longrightarrow
T_{r+r'}^{s+s'*}
(A_1,\ldots,A_r,A'_1,\ldots,A'_{r'};
 B_1,\ldots,B_s,B'_1,\ldots,B'_{s'}),
\end{aligned}
$}
\end{equation}
que satisface las propiedades evidentes de asociatividad y de compatibilidad con
los homomorfismos funtoriales, los homomorfismos de conexión y las sucesiones
espectrales. Asimismo, se tienen operaciones de simetría cuya explicitación se deja al
lector. Se tiene además una operación de contracción cada vez que
uno de los argumentos $A_i$

\src{191}
es igual a uno de los argumentos $B_j$:
\begin{equation}\tag{8.4}
\resizebox{\textwidth}{!}{$
\begin{aligned}
&T_{r}^{s*}(A_1,\ldots,A_{i-1},C,A_{i+1},\ldots,A_r;
 B_1,\ldots,B_{j-1},C,B_{j+1},\ldots,B_s)\\
&\qquad\longrightarrow
T_{r-1}^{s-1*}(A_1,\ldots,\widehat{A_i},\ldots,A_r;
 B_1,\ldots,\widehat{B_j},\ldots,B_s).
\end{aligned}
$}
\end{equation}
Además, si un argumento $A_i$ es un haz localmente libre, puede
suprimirse a condición de sustituir uno de los $A_j$ ($j\neq i$) por
$A_j\otimes A_i$, o uno de los $B_k$ por $B_k\otimes A_i'$ (donde
$A_i'=\Homsh_{\Osh_X}(A_i,\Osh_X)$), y se tiene una regla de cálculo análoga en
el caso en que uno de los argumentos $B_i$ es localmente libre. En particular, siempre
se puede suprimir un argumento que sea idéntico a $\Osh_X$. Si todos los
argumentos son localmente libres, salvo a lo sumo uno de los argumentos $B_i$, la regla
que acaba de enunciarse da un isomorfismo funtorial
\begin{equation}\tag{8.5}
  T_{r}^{s*}(A_1,\ldots,A_r;B_1,\ldots,B_s)
  =H^{*}\bigl(X,A_1'\otimes\cdots\otimes A_r'
    \otimes B_1\otimes\cdots\otimes B_s\bigr)
\end{equation}
(pues se reduce al caso en que $r=0$, $s=1$, donde esto es inmediato; también se
puede utilizar directamente la sucesión espectral de término inicial (8.2)). Para los
$\underline{T}_{r}^{s}$ se definen operaciones correspondientes a las anteriores.
Las relaciones entre las diversas operaciones así
introducidas son las mismas que para las operaciones análogas del cálculo
tensorial.

Sea $n$ la dimensión de $X$. Aplicando sucesivamente una composición
tensorial (8.3) y contracciones (8.4) sobre los argumentos repetidos, se
obtiene un acoplamiento
\begin{equation}\tag{8.6}
\resizebox{\textwidth}{!}{$
  (T_{r}^{s})^{p}(A_1,\ldots;B_1,\ldots)
  \times
  (T_{s}^{r})^{n-p}(B_1,\ldots;A_1,\ldots,A_r\otimes\Omega_X^{n})
  \longrightarrow H^{n}(X,\Omega_X^{n}).
$}
\end{equation}

\result{TEOREMA}{6} Si $X$ es una variedad proyectiva no singular, entonces
los acoplamientos (8.6) son dualidades.

Esto resulta de manera puramente formal del corolario del teorema 3. En efecto,
de este corolario resulta fácilmente que, si $K$ es un complejo de haces
algebraicos coherentes localmente libres, entonces la hipercohomología de $X$ con
respecto a $K$ está en dualidad con la hipercohomología de $X$ con respecto a
$K'\otimes\Omega_X^{n}$ mediante los acoplamientos naturales
\begin{equation}\tag{8.7}
  \mathbf R^{p}\Gamma(K)
  \times\mathbf R^{n-p}\Gamma(K'\otimes\Omega_X^{n})
  \longrightarrow\mathbf R^{n}\Gamma(\Omega_X^{n})
  =H^{n}(X,\Omega_X^{n}).
\end{equation}
Esto se ve utilizando la sucesión espectral de término inicial
$H^{p}(H^{q}(X,K))$\footnote{La rama diplomática conserva el término impreso
$H^{p}(H^{q}(X,K))$; la rama actual lo corrige según la sucesión espectral
(8.2) a $H^{p}(X,H^{q}(K))$.} y la sucesión espectral análoga para
$K'\otimes\Omega_X^{n}$. Del resultado anterior se deduce el teorema 6
utilizando la definición (8.1).

\src{192}
\subsection*{OBSERVACIONES}

\noindent 1\textsuperscript{o} Para las definiciones anteriores al teorema 6,
no era necesario que $X$ fuese no singular, pues no era
indispensable trabajar únicamente con resoluciones finitas. Pero si $X$
es singular, ya no se podrá afirmar, a priori, que los
$(\underline{T}_{r}^{s})^{p}(A_1,\ldots;B_1,\ldots)$ sean haces
coherentes, pues en el complejo de haces
\[
  \Homsh_{\Osh_X}\bigl(L(A_1)\otimes\cdots,L(B_1)\otimes\cdots\bigr)
\]
habrá una infinidad de componentes de un grado total dado.

\medskip
\noindent 2\textsuperscript{o} Se comprueba fácilmente que, en las fórmulas
(8.1), se puede sustituir uno de los $L(B_i)$ por $B_i$. Teniendo en cuenta la
proposición 3, esto demuestra que se tiene
\begin{equation}\tag{8.8}
\left\{
\begin{aligned}
  \underline{T}_{1}^{1*}(A;B)&=\Extsh_{\Osh_X}^{*}(A,B),\\
  T_{1}^{1*}(A;B)&=\operatorname{Ext}_{\Osh_X}^{*}(X;A,B).
\end{aligned}
\right.
\end{equation}
En particular, tomando en (8.6) $r=s=1$ y $A_1=\Osh_X$, se recupera el
teorema 3. La fórmula (8.8) implica también
\[
  T_{0}^{1*}(B)=H^{*}(X,B),
  \qquad
  T_{1}^{0*}(A)=\operatorname{Ext}_{\Osh_X}^{*}(X;A,\Osh_X).
\]

\medskip
\noindent 3\textsuperscript{o} Se ve en la fórmula (8.1) que, en general, los
funtores $\underline{T}_{r}^{s*}$ y $T_{r}^{s*}$ tienen componentes de
grados positivos y negativos. Utilizando la observación 2, se ve que si la
dimensión de $X$ es $n$, entonces las componentes no nulas de
$\underline{T}_{r}^{s*}$ están comprendidas entre $-(s-1)n$ y $rn$, si $s>0$,
entre $0$ y $rn$, si $s=0$, y las componentes no nulas de $T_{r}^{s*}$
están comprendidas entre $-(s-1)n$ y $(r+1)n$, si $s>0$, entre $0$ y $(r+1)n$,
si $s=0$ (e incluso, salvo error, si $r>0$, entre $-(s-1)n$ y $rn$, resp. entre
$0$ y $rn$).

\section*{BIBLIOGRAFÍA}

\noindent[1] Cartier (Pierre). --- Les groupes
$\operatorname{Ext}^{s}(A,B)$, Séminaire A. Grothendieck : Algèbre homologique,
t. 1, 1957, n\textsuperscript{o} 3.

\noindent[2] Grothendieck (Alexandre). --- Sur quelques points d'algèbre
homologique, \emph{Tohoku Math. J.}, t. 9, 1957, p. 119--183.

\noindent[3] Serre (Jean-Pierre). --- Faisceaux algébriques cohérents,
\emph{Annals of Math.}, t. 61, 1955, p. 197--278.

\noindent[4] Serre (Jean-Pierre). --- Sur la dimension homologique des anneaux
et des modules noethériens, \emph{Proc. Intern. Symp. on Alg. Number Theory}
[1955, Tokyo et Nikko]. --- Tokyo, Science Council of Japan, 1956 ;
p. 175--189.

\clearpage
\src{193}
\section*{ADENDA}

Las dificultades señaladas en la observación de la página 13 están actualmente
completamente resueltas, gracias a una extensión de los teoremas de dualidad a
variedades cualesquiera (con singularidades arbitrarias). Para esta teoría, que puede
formularse en el marco general de la « teoría de esquemas », remitimos
a los « Éléments de Géométrie algébrique » de J. Dieudonné y A. Grothendieck,
actualmente en preparación. Se encontrarán algunas indicaciones en:

\begin{quote}
Grothendieck (A.). --- Cohomology theory of algebraic varieties,
\emph{Congrès int. des Mathématiciens}, 1958, Edinburgh (de próxima aparición).
\end{quote}

\begin{flushright}
[Abril de 1959]
\end{flushright}

\begin{center}\rule{36mm}{0.5pt}\end{center}
